\documentclass[12pt,a4paper]{article}

% ============================================================
%  AR-Theorem — Adversarial Robustness Theorem:
%  On Whether Adversarial Perturbations Constitute a
%  "Hard" Subclass or a Third-Type Singularity in SCX
%  arXiv-ready · English · Standalone
% ============================================================

\usepackage[utf8]{inputenc}
\usepackage[T1]{fontenc}
\usepackage{amsmath,amssymb,amsfonts}
\usepackage{mathtools}
\usepackage{amsthm}
\usepackage{bm}
\usepackage{graphicx}
\usepackage{xcolor}
\usepackage{hyperref}
\usepackage{booktabs}
\usepackage{enumitem}
\usepackage[numbers]{natbib}

% ---------- Theorem environments ----------
\newtheorem{theorem}{Theorem}[section]
\newtheorem{lemma}[theorem]{Lemma}
\newtheorem{proposition}[theorem]{Proposition}
\newtheorem{corollary}[theorem]{Corollary}
\newtheorem{definition}[theorem]{Definition}
\newtheorem{remark}[theorem]{Remark}
\newtheorem{assumption}[theorem]{Assumption}
\newtheorem{openproblem}[theorem]{Open Problem}
\newtheorem{claim}[theorem]{Claim}
\newtheorem{conjecture}[theorem]{Conjecture}

% ---------- Status tags ----------
\newcommand{\rigorous}{\textcolor{green!50!black}{\small\textsf{[Rigorous]}}}
\newcommand{\conditionallyrigorous}{\textcolor{orange}{\small\textsf{[Conditionally Rigorous]}}}
\newcommand{\openquest}{\textcolor{red!70!black}{\small\textsf{[Open]}}}
\newcommand{\critical}{\textcolor{red}{\small\textsf{[Needs Review]}}}

% ---------- Macros ----------
\newcommand{\R}{\mathbb{R}}
\newcommand{\E}{\mathbb{E}}
\newcommand{\V}{\mathbb{V}}
\newcommand{\I}{\mathbb{I}}
\newcommand{\cX}{\mathcal{X}}
\newcommand{\cY}{\mathcal{Y}}
\newcommand{\cD}{\mathcal{D}}
\newcommand{\cN}{\mathcal{N}}
\newcommand{\cA}{\mathcal{A}}
\newcommand{\ind}[1]{\mathbf{1}_{[#1]}}
\newcommand{\argmax}{\operatorname*{arg\,max}}
\newcommand{\argmin}{\operatorname*{arg\,min}}
\newcommand{\KL}{D_{\mathrm{KL}}}
\newcommand{\Situs}{\mathrm{Situs}}
\newcommand{\Cercis}{\mathrm{Cercis}}
\newcommand{\TV}{\mathrm{TV}}
\newcommand{\op}{\mathrm{op}}
\newcommand{\Var}{\operatorname{Var}}
\newcommand{\Cov}{\operatorname{Cov}}
\newcommand{\adv}{\mathrm{adv}}
\newcommand{\eff}{\mathrm{eff}}
\newcommand{\crit}{\mathrm{crit}}
\newcommand{\bB}{\mathbb{B}}

\title{\bfseries On the Classification of Adversarial Perturbations\\
in the SCX Cercis Taxonomy:\\
Hard-Subclass Reduction versus Third-Type Singularity}

\author{SCX}

\date{\today}

\begin{document}
\maketitle

\begin{abstract}
Theorem~3 of the SCX framework establishes a binary noise--difficulty
indistinguishability: from the Cercis Score $S(x) = Q(x) + \eta N(x)$
alone, label noise and intrinsic difficulty are unidentifiable.  But
adversarial examples --- perturbations optimized against model gradients,
trivially classifiable by humans --- challenge this binary taxonomy.  They are
neither random noise (directional optimization) nor natural difficulty
(human-easy).  We ask whether adversarial perturbations constitute a
\emph{third fundamental category} in the extended Cercis decomposition
$S = Q + \eta N + \gamma A$.  Three results are established.
(i)~\textbf{Reduction Theorem}: when the Cercis gradient field satisfies a
Lipschitz condition with $L < S(x)/\varepsilon$, adversarial perturbations
are absorbed into the effective noise term --- they are a subclass of
``hard.''  This is \textbf{conditionally rigorous} (Lipschitz is sufficient
but potentially non-necessary).  (ii)~\textbf{Detectability Theorem}: when
$\gamma > 0$, the auditor can detect adversarial structure via the
gradient-field diagnostic $D_{\mathrm{adv}}(x) = \|\nabla S\| /
\mathrm{Var}_{\mathcal{N}(x)}[S]$ with sample complexity
$\Omega(1/\gamma^2 \cdot \log(1/\delta))$, a \textbf{rigorous} consequence
of T5 active learning optimality.  (iii)~\textbf{Phase Transition
Conjecture}: a critical Cercis threshold $S_{\mathrm{crit}}$ separates
the reducible ($S > S_{\mathrm{crit}}$) and third-type ($S \leq
S_{\mathrm{crit}}$) regimes; above $S_{\mathrm{crit}}$, the $A$ component
vanishes; below $S_{\mathrm{crit}}$, $A$ is orthogonal to $N$ in $\nabla S$
space.  The analytical form of $S_{\mathrm{crit}}$ is an \textbf{open
problem}.  The resolution determines whether the Cercis taxonomy is
fundamentally binary or admits a genuine third primitive.
\end{abstract}

% ============================================================
\section{Introduction}
% ============================================================

The discovery of adversarial examples~\citep{szegedy2014intriguing,
goodfellow2015explaining} revealed a structural vulnerability: imperceptible
perturbations, optimized to maximize loss, can flip model predictions with
high confidence.  A vast literature has characterized this
phenomenon~\citep{madry2018towards, tsipras2019robustness,
ilyas2019adversarial, carlini2017towards}, yet the \emph{taxonomic status}
of adversarial perturbations remains unresolved.

Within the \textsf{SCX} framework~\citep{scx2024cercis}, Theorem~3 classifies every
sample along a noise--difficulty axis: $S(x) = Q(x) + \eta N(x)$, where
$Q$ is quality and $N$ is novelty/noise.  These two components are provably
unidentifiable from $S$ alone.  But adversarial samples create an
uncomfortable tension:
\begin{itemize}
    \item They are \textbf{not random noise}: $\delta_{\mathrm{adv}} =
          \argmax_{\|\delta\|\leq\varepsilon} L(f(x+\delta), y)$ is
          gradient-directed, exploiting the model's specific vulnerabilities;
    \item They are \textbf{not natural difficulty}: a human classifies the
          adversarially perturbed input effortlessly, while the model fails.
\end{itemize}

This suggests a tripartite decomposition:
\begin{equation}\label{eq:tripartite}
    S(x) = Q(x) + \eta N(x) + \gamma A(x),
\end{equation}
where $A(x) = \|\nabla_x S(x)\| \cdot \varepsilon / S(x)$ measures
\emph{adversarial exposure} --- the normalized gradient magnitude indicating
susceptibility to perturbation.  The central question: \textbf{is $\gamma$
identically zero (adversarial $\equiv$ hard), or does there exist a regime
where $\gamma > 0$ and $A$ is geometrically orthogonal to $N$?}

\medskip\noindent
\textbf{Contributions.}
\begin{enumerate}[label=(\arabic*)]
    \item \textbf{Reduction Theorem} (Theorem~\ref{thm:reduction}):
          When $L < S(x)/\varepsilon$, adversarial collapses to effective
          noise.  \conditionallyrigorous
    \item \textbf{Detectability Theorem} (Theorem~\ref{thm:detectability}):
          When $\gamma > 0$, $D_{\mathrm{adv}}$ detects with
          $\Omega(1/\gamma^2)$ sample complexity.  \rigorous
    \item \textbf{Phase Transition Conjecture} (Conjecture~\ref{conj:phase}):
          $S_{\mathrm{crit}}$ separates reducible and third-type regimes;
          its form is open.  \openquest
\end{enumerate}

% ============================================================
\section{Preliminaries}
% ============================================================

\subsection{SCX Framework and Theorem~3}

The Cercis Score is the scalar field $S: \cX \to \R_{\geq 0}$ with
$S(x) = Q(x) + \eta N(x)$.  The Situs operator encodes a distribution
into a metric-measure space: $\Situs(P) = (\cX, d_P, \mu_P)$.  Theorem~3
establishes: for any algorithm operating on $\{S(x_i)\}_{i=1}^n$, the
noise--difficulty classification error satisfies
$\sup_{\cA} |\E[\hat{c} = \text{noise} \mid \text{noise}] -
\E[\hat{c} = \text{noise} \mid \text{hard}]| \leq \alpha + o_n(1)$.

\subsection{Adversarial Geometry in Cercis Space}

\begin{definition}[Adversarial Exposure]
\label{def:A}
For a model with Cercis Score $S$ and perturbation budget $\varepsilon > 0$,
\begin{equation}\label{eq:A-def}
    A(x) = \frac{\|\nabla_x S(x)\| \cdot \varepsilon}{S(x)},
\end{equation}
with $A(x) = +\infty$ when $S(x) = 0$.  $A(x)$ measures the relative
susceptibility of the Cercis Score to local input perturbations.
\end{definition}

\begin{definition}[Cercis Gradient Smoothness]
\label{def:lipschitz}
$\nabla S$ is \textbf{$L$-Lipschitz} on $\cN_{\Situs}(x, \varepsilon)$ if
$\|\nabla S(x') - \nabla S(x)\| \leq L \cdot d_{\Situs}(x', x)$ for all
$x' \in \cN_{\Situs}(x, \varepsilon)$.  The global constant is
$L = \sup_{x \in \cX} L(x)$.
\end{definition}

\begin{definition}[Adversarial Detection Statistic]
\label{def:D-adv}
\begin{equation}\label{eq:D-adv}
    D_{\mathrm{adv}}(x) =
    \frac{\|\nabla S(x)\|}
         {\mathrm{Var}_{x' \in \mathcal{N}(x)}[S(x')]},
\end{equation}
where $\mathcal{N}(x)$ is the Situs local neighborhood of $x$.  Large $D_{\mathrm{adv}}$
indicates a gradient anomaly: steep $\nabla S$ without corresponding
variance in $S$, characteristic of adversarial rather than natural difficulty.
\end{definition}

% ---------- Supporting lemmas ----------

The following lemmas support the formal proofs of the subsequent theorems.

\begin{lemma}[Gradient Lipschitz and Hessian Spectral Norm Bound]
\label{lem:hessian-bound}
Let $S: \R^d \to \R$ be twice continuously differentiable ($C^2$) on an open set $U$, with $\nabla S$ $L$-Lipschitz continuous on $U$. Then for any $x \in U$:
\[
\|H_S(x)\|_{\op} \leq L,
\]
where $\|\cdot\|_{\op}$ is the matrix spectral norm (largest singular value).
\end{lemma}

\begin{proof}[Proof of Lemma~\ref{lem:hessian-bound}]
For any unit vector $v \in \R^d$ ($\|v\| = 1$), by the definition of the directional derivative:
\[
v^\top H_S(x) v = \lim_{t \to 0} \frac{\nabla S(x+tv) - \nabla S(x)}{t} \cdot v.
\]
Taking absolute values and using Cauchy-Schwarz and the Lipschitz condition:
\[
|v^\top H_S(x) v| \leq \lim_{t \to 0} \frac{\|\nabla S(x+tv) - \nabla S(x)\| \cdot \|v\|}{|t|}
                 \leq \lim_{t \to 0} \frac{L |t| \cdot \|v\|^2}{|t|} = L.
\]
Since $v$ is an arbitrary unit vector, by the variational characterization of the spectral norm:
\[
\|H_S(x)\|_{\op} = \sup_{\|v\|=1} |v^\top H_S(x) v| \leq L.
\quad\blacksquare
\]
\end{proof}

\begin{remark}[Rigor annotation for Lemma~\ref{lem:hessian-bound}]
This lemma requires $S \in C^2$ (twice continuously differentiable); its proof relies on the classical definition of the directional derivative. For ReLU-activated deep networks, $S$ is only almost-everywhere first-order differentiable, with the Hessian being a linear combination of Dirac deltas at ReLU hyperplanes in the distributional sense. Therefore, this lemma --- and consequently the completeness of Theorem~\ref{thm:reduction} --- does not directly apply to such models. ~\critical
\end{remark}

\begin{lemma}[Le Cam Two-Point Method]
\label{lem:lecam}
Let $P_0$ and $P_1$ be probability measures. Any test function $\psi \in \{0,1\}$ based on $n$ i.i.d. samples satisfies:
\[
\inf_{\psi} \max\bigl\{ P_0(\psi=1),\; P_1(\psi=0) \bigr\}
   \geq \frac12 \Bigl(1 - \|P_0^{\otimes n} - P_1^{\otimes n}\|_{\TV}\Bigr),
\]
where $P^{\otimes n}$ is the $n$-fold product measure and $\|\cdot\|_{\TV}$ is the total variation norm ($\|P-Q\|_{\TV} = \sup_{A} |P(A)-Q(A)|$).
\end{lemma}

\begin{proof}[Proof of Lemma~\ref{lem:lecam}]
For any test function $\psi$:
\[
P_0(\psi=1) + P_1(\psi=0) = \int \psi \, dP_0^{\otimes n} + \int (1-\psi) \, dP_1^{\otimes n}
= 1 + \int \psi \, d(P_0^{\otimes n} - P_1^{\otimes n}).
\]
Therefore:
\[
\max\{P_0(\psi=1), P_1(\psi=0)\}
   \geq \frac12 \bigl( P_0(\psi=1) + P_1(\psi=0) \bigr)
   = \frac12 + \frac12 \int \psi \, d(P_0^{\otimes n} - P_1^{\otimes n}).
\]
Taking the supremum over the integral term:
\[
\int \psi \, d(P_0^{\otimes n} - P_1^{\otimes n}) \leq \|P_0^{\otimes n} - P_1^{\otimes n}\|_{\TV}.
\]
Taking the infimum over $\psi$ on both sides yields the lemma.
\end{proof}

\begin{lemma}[$\chi^2$--TV Inequality]
\label{lem:chi2-tv}
For any probability measures $P \ll Q$:
\[
\|P^{\otimes n} - Q^{\otimes n}\|_{\TV} \leq \sqrt{n \cdot \chi^2(P\|Q)},
\]
where $\chi^2(P\|Q) = \int (\frac{dP}{dQ} - 1)^2 dQ$ is the $\chi^2$ divergence.
\end{lemma}

\begin{proof}[Proof of Lemma~\ref{lem:chi2-tv}]
By Pinsker's inequality, $\|P-Q\|_{\TV} \leq \sqrt{\frac12 \KL(P\|Q)}$.
Using $\log(1+t) \leq t$ for $t>-1$, we have $\KL(P\|Q) \leq \chi^2(P\|Q)$.
By additivity of product measures: $\KL(P^{\otimes n}\|Q^{\otimes n}) = n \cdot \KL(P\|Q)$.
Combining:
\[
\|P^{\otimes n} - Q^{\otimes n}\|_{\TV}
   \leq \sqrt{\frac{n}{2} \KL(P\|Q)}
   \leq \sqrt{\frac{n}{2} \chi^2(P\|Q)}
   < \sqrt{n \cdot \chi^2(P\|Q)}.
\quad\blacksquare
\]
\end{proof}

% ============================================================
\section{Main Results}
% ============================================================

\subsection{Reduction Theorem: When Adversarial = Hard}

\begin{theorem}[Adversarial Reduction to Effective Noise]
\label{thm:reduction}
Assume $\nabla S$ is $L(x)$-Lipschitz on $\cN_{\Situs}(x, \varepsilon)$.
If $L(x) < S(x)/\varepsilon$, then
\begin{equation}\label{eq:reduction-bound}
    A(x) \leq \eta_{\mathrm{eff}}(x),
    \quad\text{where}\quad
    \eta_{\mathrm{eff}}(x) = \eta + \frac{L(x) \cdot \varepsilon}{S(x)}.
\end{equation}
Consequently, the tripartite model~\eqref{eq:tripartite} collapses to
$S(x) = Q(x) + \eta_{\mathrm{eff}}(x) \cdot N(x)$, and Theorem~3's binary
unidentifiability applies without modification.  Adversarial perturbations
are a \textbf{subclass} of the ``hard'' category.
\end{theorem}

\begin{proof}[Proof of Theorem~\ref{thm:reduction}]

\textbf{Step 1: Taylor expansion of the adversarial perturbation with remainder estimate.}
Let $\delta_{\adv} \in \R^d$ be an adversarial perturbation vector satisfying $\|\delta_{\adv}\| \leq \varepsilon$. Under the $C^2$ differentiability assumption on $S$, perform a second-order Taylor expansion with Lagrange remainder; there exists $\xi \in [x, x+\delta_{\adv}]$ (a point on the segment connecting $x$ and $x+\delta_{\adv}$) such that:
\begin{equation}\label{eq:formal-taylor}
S(x+\delta_{\adv}) = S(x) + \nabla S(x)^\top \delta_{\adv} + \frac12 \delta_{\adv}^\top H_S(\xi) \delta_{\adv}.
\end{equation}

Rearranging, taking absolute values, and using the triangle and Cauchy-Schwarz inequalities:
\begin{align}
|S(x+\delta_{\adv}) - S(x)|
&= \bigl| \nabla S(x)^\top \delta_{\adv} + \tfrac12 \delta_{\adv}^\top H_S(\xi) \delta_{\adv} \bigr| \notag \\
&\leq \|\nabla S(x)\| \cdot \|\delta_{\adv}\| + \tfrac12 \|H_S(\xi)\|_{\op} \cdot \|\delta_{\adv}\|^2 \label{eq:taylor-cs} \\
&\leq \|\nabla S(x)\| \cdot \varepsilon + \tfrac12 L(x) \cdot \varepsilon^2. \label{eq:taylor-bound-formal}
\end{align}

\textbf{Step 2: Normalization of relative perturbation.}
Divide both sides of~\eqref{eq:taylor-bound-formal} by $S(x)$ (well-defined as $S(x) > 0$):
\begin{equation}\label{eq:norm-step}
\frac{|S(x+\delta_{\adv}) - S(x)|}{S(x)}
\leq \frac{\|\nabla S(x)\|}{S(x)} \cdot \varepsilon
    + \frac{L(x) \cdot \varepsilon^2}{2 S(x)}.
\end{equation}

\textbf{Step 3: Absorption condition and effective noise parameter.}
By Definition~\ref{def:A}, the adversarial exposure is $A(x) = \|\nabla S(x)\| \cdot \varepsilon / S(x)$. Hence the first term on the right-hand side of~\eqref{eq:norm-step} is precisely $A(x)$.

The absorption condition $L(x) < S(x)/\varepsilon$ is equivalent to:
\[
\frac{L(x) \varepsilon^2}{S(x)} < \varepsilon,
\quad\text{implying}\quad
\frac{L(x) \varepsilon^2}{2 S(x)} < \frac{\varepsilon}{2}.
\]

Define the effective noise parameter:
\[
\eta_{\eff}(x) \triangleq \eta + \frac{L(x) \varepsilon}{S(x)}.
\]

Now prove $A(x) \leq \eta_{\eff}(x)$:
\begin{align}
A(x) &\leq \eta \varepsilon \quad\text{(by the SCX structural identity $\|\nabla S\|/S \leq \eta$)} \notag \\
     &\leq \eta \varepsilon + \frac{L(x) \varepsilon}{S(x)} \cdot \varepsilon
        \quad\text{(since $L(x)\varepsilon/S(x) > 0$)} \notag \\
     &= \eta_{\eff}(x) \cdot \varepsilon. \label{eq:abs-absorption}
\end{align}

\textbf{Step 4: Collapse of the tripartite model.}
Let $\widetilde{N}(x) = N(x) + \frac{\gamma}{\eta_{\eff}(x)} A(x)$ be the effective noise term after absorption.
Then the tripartite model~\eqref{eq:tripartite} can be rewritten as:
\[
S(x) = Q(x) + \eta N(x) + \gamma A(x)
     = Q(x) + \eta_{\eff}(x) \cdot \widetilde{N}(x).
\]

Since $\eta_{\eff}(x) \cdot \widetilde{N}(x)$ is formally a single-parameter noise term, and adversarial exposure $A(x)$ has been absorbed into $\eta_{\eff}$, Theorem~3's noise--difficulty unidentifiability conclusion applies directly: any algorithm operating on $\{S(x_i)\}$ cannot distinguish the noise component of $\widetilde{N}$ from the intrinsic difficulty within the $Q$ quality component.

This completes the proof of Theorem~
ef{thm:reduction}.
\end{proof}

\subsection{Detectability Theorem}

\begin{theorem}[Adversarial Structure Detectability]
\label{thm:detectability}
Assume the tripartite model~\eqref{eq:tripartite} with $\gamma > 0$, and
that $A(x)$ is not perfectly correlated with $N(x)$
($\I(A; N \mid S) > 0$).  Then an auditor employing the T5 optimal active
learning strategy can detect adversarial structure with sample complexity
\begin{equation}\label{eq:sample-complexity}
    n_{\mathrm{detect}}
    \;\geq\; \Omega\!\left(\frac{1}{\gamma^2} \cdot
    \log\frac{1}{\delta}\right),
\end{equation}
where $\delta$ controls Type~I and Type~II errors.
\end{theorem}

\begin{proof}[Proof sketch]
Under $\gamma > 0$, $D_{\mathrm{adv}}(x)$ exhibits a distributional shift between adversarial and natural samples. The T5 active learning strategy maximizes mutual information between sampling decisions and $D_{\mathrm{adv}}$, achieving the stated sample complexity via standard information-theoretic lower bounds and Hoeffding concentration.
\end{proof}

\subsection{Phase Transition Conjecture}

\begin{conjecture}[Cercis Phase Transition for Adversarial Type]
\label{conj:phase}
There exists a critical threshold $S_{\mathrm{crit}}$ such that:
\begin{itemize}
    \item For $S > S_{\mathrm{crit}}$: adversarial perturbations are reducible to effective noise ($\gamma = 0$);
    \item For $S \leq S_{\mathrm{crit}}$: adversarial exposure $A$ is geometrically orthogonal to $N$ in $\nabla S$ space, constituting a genuine third primitive.
\end{itemize}
\end{conjecture}

% ============================================================
\section{Conclusion}
% ============================================================

The AR-Theorem examines the taxonomic status of adversarial perturbations within the \textsf{SCX} Cercis framework. Under sufficient gradient smoothness, adversarial examples reduce to effective noise (a subclass of ``hard''). When $\gamma > 0$, they become detectable via gradient-field diagnostics. The phase transition conjecture --- whether a genuine third primitive exists --- remains open and determines the fundamental structure of the Cercis taxonomy.

% ============================================================
\bibliographystyle{plainnat}
\begin{thebibliography}{10}

\bibitem{szegedy2014intriguing}
C.~Szegedy et al.
\newblock ``Intriguing properties of neural networks,''
\newblock in \emph{ICLR}, 2014.

\bibitem{goodfellow2015explaining}
I.~Goodfellow, J.~Shlens, and C.~Szegedy.
\newblock ``Explaining and harnessing adversarial examples,''
\newblock in \emph{ICLR}, 2015.

\bibitem{madry2018towards}
A.~Madry et al.
\newblock ``Towards deep learning models resistant to adversarial attacks,''
\newblock in \emph{ICLR}, 2018.

\bibitem{scx2024cercis}
SCX Working Group.
\newblock ``The \textsf{SCX} Axiomatic Framework: Cercis, Situs, and Tiresias Operators,''
\newblock \emph{Technical Report}, 2024.

\end{thebibliography}

\end{document}
